\section{Epistemic logic}
\label{sec:modal-temporal-epistemic-03}

Fix a nonempty set \(I\) of agents.  The language adds \(K_iA\) for each
\(i\in I\).  An epistemic frame is
\[
  (W,(\sim_i)_{i\in I}),
\]
where \(W\ne\varnothing\) and each
\(\sim_i\subseteq W\times W\).  An epistemic structure is
\(\mathcal M=(W,(\sim_i)_{i\in I},V)\), where
\(V:\Prop\to\Pow(W)\), with the classical connective clauses and
\[
  \mathcal M,w\entails K_iA
  \quad\Longleftrightarrow\quad
  \forall v\,(w\sim_i v\to\mathcal M,v\entails A).
\]
For knowledge, each \(\sim_i\) is normally required to be an equivalence
relation.  The multi-agent calculus used below has all uniform substitution
instances of classically valid propositional formulas and, for each
\(i\in I\), the
schemas
\[
\begin{aligned}
  K_i(A\to B)&\to(K_iA\to K_iB),&
  K_iA&\to A,\\
  K_iA&\to K_iK_iA,&
  \neg K_iA&\to K_i\neg K_iA.
\end{aligned}
\]
For a premise set \(\Gamma\), a derivation is a finite sequence in which each
line is a member of \(\Gamma\), an instance of one of the displayed schemas,
follows by modus ponens, or has the form \(K_iA\) where \(A\) has a separate
no-premise derivation.  Write \(\Gamma\derives_{\Sfive_I}A\) when such a
derivation ends in \(A\).  This theorem-only restriction on necessitation is
part of the calculus.  The subscript \(I\) records that the construction is the
componentwise multi-agent extension of the previously defined \(\Sfive\)
calculus; it is needed to distinguish the agent-indexed language and adds no
interaction axioms between distinct agents.

